% GENERATED FILE. Edit canonical lessons, then run npm run generate:books.
% canonical-source-hash: fnv1a64:5188d177b18493ab
% canonical-lessons: IT-C511-segnare, IT-C511-pareggio, IT-C511-pareggiare, IT-C511-classifica, IT-C511-eliminare

\chapter{Il risultato}
\label{ch:it-il-risultato}
\hlchaptermodality{511}

\begin{chapteropening}
\textbf{By the end of this chapter:} \emph{I can talk about results: goals, draws and the league table.}

Complete the last lesson of chapter 511: I can talk about results: goals, draws and the league table.
\end{chapteropening}

\section[segnare]{segnare — to score; to mark, to note down}

\label{lesson:IT-C511-segnare}



\begin{quote}
Before the new one: say the Italian for first-team player, then the Italian for to save.
\end{quote}

\subsection*{You'll want to know: segnare}
\textbf{segnare} — \textquotedblleft{}to score; to mark, to note down\textquotedblright{}.

\begin{quote}
Ha segnato due gol nel primo tempo.
\end{quote}

\begin{quote}
\emph{He scored two goals in the first half.}
\end{quote}

Also to note down: \emph{segnare un numero}.

\begin{cousinweb}
\emph{segnare} is inherited from Latin \emph{signāre} \textquotedblleft{}to mark, to seal\textquotedblright{} (EtymDB), from \emph{signum} \textquotedblleft{}a mark\textquotedblright{}.
\end{cousinweb}

\begin{grammarlens}[title={What you've built}]
One more word for telling, explaining and describing at B1.
\end{grammarlens}

\subsection*{Your turn}
\begin{itemize}
\raggedright
  \item \emph{Say it:} \emph{segnare}
  \item \emph{Say it:} \emph{segnare}, once more
  \item \emph{Say it:} a sentence of your own with \emph{segnare}
  \item \emph{Recall:} say the Italian for first-team player, then the Italian for to save, then say \emph{segnare} again
\end{itemize}

\begin{tcolorbox}[breakable,colback=teal!4,colframe=teal!35!black,title={Before you move on}]
What does segnare mean? (\textquotedblleft{}To score; to mark, to note down\textquotedblright{}.) Say it once more.
\end{tcolorbox}
\section[il pareggio]{il pareggio — draw, tie}

\label{lesson:IT-C511-pareggio}



\begin{quote}
Before the new one: say the Italian for to save, then the Italian for to score.
\end{quote}

\subsection*{You'll want to know: il pareggio}
\textbf{il pareggio} — \textquotedblleft{}draw, tie\textquotedblright{}.

\begin{quote}
La partita è finita con un pareggio.
\end{quote}

\begin{quote}
\emph{The match ended in a draw.}
\end{quote}

Built on \emph{pari} \textquotedblleft{}equal\textquotedblright{}.

\begin{cousinweb}
\emph{pareggio} comes from \emph{pareggiare} \textquotedblleft{}to make even\textquotedblright{} (EtymDB), built on \emph{pari} \textquotedblleft{}equal\textquotedblright{}, from Latin \emph{pār, paris} (Lewis \& Short).
\end{cousinweb}

\begin{grammarlens}[title={What you've built}]
One more word for telling, explaining and describing at B1.
\end{grammarlens}

\subsection*{Your turn}
\begin{itemize}
\raggedright
  \item \emph{Say it:} \emph{il pareggio}
  \item \emph{Say it:} \emph{il pareggio}, once more
  \item \emph{Say it:} a sentence of your own with \emph{il pareggio}
  \item \emph{Recall:} say the Italian for to save, then the Italian for to score, then say \emph{il pareggio} again
\end{itemize}

\begin{tcolorbox}[breakable,colback=teal!4,colframe=teal!35!black,title={Before you move on}]
What does il pareggio mean? (\textquotedblleft{}Draw, tie\textquotedblright{}.) Say it once more.
\end{tcolorbox}
\section[pareggiare]{pareggiare — to draw, to equalize}

\label{lesson:IT-C511-pareggiare}



\begin{quote}
Before the new one: say the Italian for to score, then the Italian for draw.
\end{quote}

\subsection*{You'll want to know: pareggiare}
\textbf{pareggiare} — \textquotedblleft{}to draw, to equalize\textquotedblright{}.

\begin{quote}
Abbiamo pareggiato uno a uno.
\end{quote}

\begin{quote}
\emph{We drew one all.}
\end{quote}

Also to even out: \emph{pareggiare i conti}.

\begin{cousinweb}
\emph{pareggiare} is built on \emph{pari} \textquotedblleft{}equal, even\textquotedblright{}, inherited from Latin \emph{pār, paris} \textquotedblleft{}equal\textquotedblright{} (Lewis \& Short).
\end{cousinweb}

\begin{grammarlens}[title={What you've built}]
One more word for telling, explaining and describing at B1.
\end{grammarlens}

\subsection*{Your turn}
\begin{itemize}
\raggedright
  \item \emph{Say it:} \emph{pareggiare}
  \item \emph{Say it:} \emph{pareggiare}, once more
  \item \emph{Say it:} a sentence of your own with \emph{pareggiare}
  \item \emph{Recall:} say the Italian for to score, then the Italian for draw, then say \emph{pareggiare} again
\end{itemize}

\begin{tcolorbox}[breakable,colback=teal!4,colframe=teal!35!black,title={Before you move on}]
What does pareggiare mean? (\textquotedblleft{}To draw, to equalize\textquotedblright{}.) Say it once more.
\end{tcolorbox}
\section[la classifica]{la classifica — league table, ranking; chart}

\label{lesson:IT-C511-classifica}



\begin{quote}
Before the new one: say the Italian for draw, then the Italian for to draw.
\end{quote}

\subsection*{You'll want to know: la classifica}
\textbf{la classifica} — \textquotedblleft{}league table, ranking; chart\textquotedblright{}.

\begin{quote}
Siamo terzi in classifica.
\end{quote}

\begin{quote}
\emph{We are third in the league.}
\end{quote}

For music too: \emph{la classifica dei dischi}.

\begin{cousinweb}
\emph{classifica} comes from \emph{classificare} \textquotedblleft{}to classify\textquotedblright{}, built on Latin \emph{classis} \textquotedblleft{}a class, a division\textquotedblright{} (Lewis \& Short) and \emph{facere} \textquotedblleft{}to make\textquotedblright{}.
\end{cousinweb}

\begin{grammarlens}[title={What you've built}]
One more word for telling, explaining and describing at B1.
\end{grammarlens}

\subsection*{Your turn}
\begin{itemize}
\raggedright
  \item \emph{Say it:} \emph{la classifica}
  \item \emph{Say it:} \emph{la classifica}, once more
  \item \emph{Say it:} a sentence of your own with \emph{la classifica}
  \item \emph{Recall:} say the Italian for draw, then the Italian for to draw, then say \emph{la classifica} again
\end{itemize}

\begin{tcolorbox}[breakable,colback=teal!4,colframe=teal!35!black,title={Before you move on}]
What does la classifica mean? (\textquotedblleft{}League table, ranking; chart\textquotedblright{}.) Say it once more.
\end{tcolorbox}
\section[eliminare]{eliminare — to knock out, to eliminate}

\label{lesson:IT-C511-eliminare}



\begin{quote}
Before the new one: say the Italian for to draw, then the Italian for league table.
\end{quote}

\subsection*{You'll want to know: eliminare}
\textbf{eliminare} — \textquotedblleft{}to knock out, to eliminate\textquotedblright{}.

\begin{quote}
La nostra squadra è stata eliminata ai quarti.
\end{quote}

\begin{quote}
\emph{Our team was knocked out in the quarter-finals.}
\end{quote}

\emph{Essere eliminati} — to be knocked out.

\begin{cousinweb}
\emph{eliminare} is a learned borrowing from Latin \emph{ēlīmināre} \textquotedblleft{}to put out of doors\textquotedblright{} (Lewis \& Short), from \emph{ē-} and \emph{līmen} \textquotedblleft{}threshold\textquotedblright{}.
\end{cousinweb}

\begin{grammarlens}[title={What you've built}]
One more word for telling, explaining and describing at B1.
\end{grammarlens}

\subsection*{Your turn}
\begin{itemize}
\raggedright
  \item \emph{Say it:} \emph{eliminare}
  \item \emph{Say it:} \emph{eliminare}, once more
  \item \emph{Say it:} a sentence of your own with \emph{eliminare}
  \item \emph{Recall:} say the Italian for to draw, then the Italian for league table, then say \emph{eliminare} again
\end{itemize}

\begin{tcolorbox}[breakable,colback=teal!4,colframe=teal!35!black,title={Before you move on}]
What does eliminare mean? (\textquotedblleft{}To knock out, to eliminate\textquotedblright{}.) Say it once more.
\end{tcolorbox}
